\documentclass[10pt,a4paper]{amsart}
\usepackage[margin=19mm]{geometry}
\usepackage{amsmath,amssymb,mathtools}
\usepackage[scr=rsfs,cal=euler]{mathalfa}
\usepackage{xfrac}
\usepackage[unicode,hidelinks,bookmarks=false]{hyperref}
\newcommand{\ie}{i.e.\ }
\newcommand{\cf}{cf.\ }
\newcommand{\resp}{respectively\ }
\newcommand{\Subsection}{\subsection}
\providecommand{\nobreakdash}{\nobreak\hbox{-}}
\setlength{\emergencystretch}{2em}
\multlinegap0pt
\usepackage{bm}
\usepackage{bbm}
\usepackage{dsfont}
\usepackage{xspace}
\usepackage{cancel}
\usepackage{mathtools}
\usepackage{amsthm}
\makeatletter
\@ifundefined{theorem}{\newtheorem{theorem}{Theorem}}{}
\@ifundefined{lemma}{\newtheorem{lemma}[theorem]{Lemma}}{}
\makeatother
\newcommand{\calL}{\mathcal{L}}
\newcommand{\calM}{\mathcal{M}}
\newcommand{\impl}{\rightarrow} % logical implication
\newcommand{\lfp}{\mathrm{lfp}} % least-fixed point
\newcommand{\modelsa}{\models_\mathrm{a}} % models abstractly
\newcommand{\mybar}[1]{\overline{#1}}
\title[An abstract fixed-point theorem for Horn formula equations]{An abstract fixed-point theorem for Horn formula equations}
\author{Stefan Hetzl, Johannes Kloibhofer}
\date{Source: arXiv:2511.08162v1 (CC BY 4.0)}
\begin{document}
\maketitle

\noindent\emph{Source and licence.} An abstract fixed-point theorem for Horn formula equations. Version arXiv:2511.08162v1 (CC BY 4.0), distributed under the Creative Commons Attribution 4.0 International licence, as recorded in the arXiv licence metadata for this exact version and shown on the abstract page https://arxiv.org/abs/2511.08162v1. Full rights notice: licenses/arxiv-2511.08162-CC-BY-4.0.txt.\par\smallskip

\noindent\textit{Editorial scope.} Counted unit: the Lemma whose source label is \texttt{lem.abstractNonnengart} in the article (source0.tex lines 1134--1140, source label \texttt{lem.abstractNonnengart}), delivered together with the article's own complete original proof of it (source0.tex lines 1141--1183, one \texttt{proof} environment of 43 lines). No mathematics is written, repaired or re-derived: statement and proof are the article's own text line for line. Required context retained uncounted: lem.ImplPositiv (lines 356--365). Every definition, display and label of those lines is retained unchanged. Omitted from this package and not redistributed: 1--355, 366--1133, 1184--2376. Nothing inside a retained excerpt is omitted or altered: every retained body is delivered exactly as the article's lines are written, so no omission receipt of any kind accompanies this package. Citations: the retained excerpts carry no citation command at all, so no bibliography entry of the article is transcribed and no reference is redistributed. Reference adaptation: none was needed and none was made; every reference inside the delivered unit resolves inside it, so no unresolved reference or citation is produced. Printed slips kept literal: no printed word, symbol, hypothesis or number of the counted statement or of its proof was altered. Layout and class: the article's own class is \texttt{acmart} with packages including inputenc, amsmath, ebproof, amsthm, fixme; the wrapper is the lane's pinned \texttt{amsart} editorial wrapper at 10pt on A4, which restates the theorem-like environments the retained text needs and adds the article's own macro definitions that the retained text needs (\texttt{\textbackslash calL}, \texttt{\textbackslash calM}, \texttt{\textbackslash impl}, \texttt{\textbackslash lfp}, \texttt{\textbackslash modelsa}, \texttt{\textbackslash mybar}), with extra wrapper packages (bm, bbm, dsfont, xspace, cancel, mathtools). The delivered numbering is the unit's own. Assets: no retained span contains a figure environment, a graphics-inclusion command, a PDF-page inclusion, a table or a TikZ picture; no image file is copied into this package, the package declares no asset dependency and no image file is redistributed with it. Licence: version arXiv:2511.08162v1 of this article is distributed under the Creative Commons Attribution 4.0 International licence, as recorded in the arXiv licence metadata for this exact version and shown on the abstract page https://arxiv.org/abs/2511.08162v1. A full rights notice is delivered at licenses/arxiv-2511.08162-CC-BY-4.0.txt. Characters: every retained excerpt is pure ASCII, so the delivered engine, XeLaTeX with its default Latin Modern text font, reports no missing character.
\begin{lemma}\label{lem.ImplPositiv}
	Let $\calL$ be a language and $\calM$ be an $\calL$-structure. Let $R$ be a predicate variable and let $S$ and $S'$ be relations in $\calM$ with the same arity as $R$. Let $\varphi$ be a first-order formula in $\calL \cup \{R\}$, such that $R$ occurs only positively in $\varphi$. Then
	\begin{align*}
		\calM \models S(\mybar{x}) \rightarrow S'(\mybar{x}) \quad \Rightarrow \quad \calM \models \varphi[R \backslash S] \rightarrow \varphi[R \backslash S'].
	\end{align*}
	Conversely, if $R$ occurs only negatively in $\varphi$, then
	\begin{align*}
		\calM \models S(\mybar{x}) \rightarrow S'(\mybar{x}) \quad \Rightarrow \quad \calM \models \varphi[R \backslash S'] \rightarrow \varphi[R \backslash S].
	\end{align*}
\end{lemma}

\begin{lemma}\label{lem.abstractNonnengart}
	Let $X_1,...,X_n$ be predicate variables. Let $\beta(X_1,...,X_n)$ be a first-order formula and $\Phi = (\alpha_i(X_1,...,X_n,\mybar{x_i},))_{i=1}^n$ be an $n$-tuple of first-order formulas such that $|\mybar{x_i}|$ equals the arity of $X_i$ for $i \in \{1,...,n\}$. \footnote{Note that there could also be free variables in $\alpha_1,...,\alpha_n,\beta$, that are not stated explicitly.} If $X_1,...,X_n$ occur only positively in $\alpha_1,...,\alpha_n$ and only negatively in $\beta$, then 
	\begin{enumerate}
		\item $\modelsa \exists \mybar{X} \left(\bigwedge_{i=1}^n \forall \mybar{x_i} \left(\alpha_i(X_1,...,X_n,\mybar{x_i}) \impl X_i(\mybar{x_i})\right) \wedge \beta(X_1,...,X_n) \right) \leftrightarrow \beta([\lfp_{X_1} ~\Phi],...,[\lfp_{X_n} ~\Phi])$,
		\item If  $(\calM,G) \modelsa \bigwedge_{i=1}^n \forall \mybar{x_i} \left(\alpha_i( R_1,...,R_n,\mybar{x_i}) \impl R_i(\mybar{x_i})\right)$ for some model abstraction $(\calM,G)$ and abstract relations $R_1,...,R_n$, then $(\calM,G) \modelsa \bigwedge_{i=1}^n \forall \mybar{x_i} ([\lfp_{X_i} ~\Phi](\mybar{x_i}) \impl R_i(\mybar{x_i}))$.
	\end{enumerate}
\end{lemma}

\begin{proof}
	1. $``\rightarrow"$: Let $(\calM,G)$ be a model abstraction and $\theta$ an environment such that
	\[
	(\calM,G),\theta \modelsa \exists \mybar{X} \bigwedge_{i=1}^n \forall \mybar{x_i} \left(\alpha_i(X_1,...,X_n,\mybar{x_i}) \impl X_i(\mybar{x_i})\right) \wedge \beta(X_1,...,X_n).
	\]
	Hence there exist abstract relations $S_1,...,S_n$ and an environment $\theta' := \theta[X_1 := \gamma_{k_1}(S_1),...,X_n := \gamma_{k_n}(S_n)]$  such that 
	\begin{align}\label{abstractNonnengartproof}
		(\calM,G),\theta' \modelsa \bigwedge_{i=1}^n \forall \mybar{x_i} \left(\alpha_i(X_1,...,X_n,\mybar{x_i}) \impl X_i(\mybar{x_i})\right) \wedge \beta(X_1,...,X_n).
	\end{align}
Then $(F_{\Phi})_i(\gamma_{k_1}(S_1),...,\gamma_{k_n}(S_n)) \subseteq \gamma_{k_i}(S_i)$ for all $i \in \{1,...,n\}$ and as $\alpha_k,\gamma_k$ form a Galois connection for every $k$ we have $\alpha_{k_i} \circ (F_{\Phi})_i(\gamma_{k_1}(S_1),...,\gamma_{k_n}(S_n)) \sqsubseteq S_i$ for all $i \in \{1,...,n\}$. Thus $(S_1,...,S_n)$ is a fixed point of $F_{\Phi}^{\#}$. Using the monotonicity of $\gamma_k$ for every $k$, this implies
\begin{align}\label{abstactNonnengartproofPart2}
	(\calM,G),\theta' \modelsa \bigwedge_{i=1}^n \forall \mybar{x_i} ([\lfp_{X_i} ~\Phi](\mybar{x_i}) \impl X_i(\mybar{x_i})).
\end{align}

	Note that the FO[LFP]-formulas are well defined, as $X_1,...,X_n$ occur only positively in $\alpha_1,...,\alpha_n$. From (\ref{abstractNonnengartproof}) we also obtain $(\calM,G),\theta' \modelsa \beta(X_1,...,X_n)$ and as $X_1,...,X_n$ occur only negatively in $\beta$ we may use Lemma \ref{lem.ImplPositiv} to obtain
	\[
	(\calM,G),\theta' \modelsa \beta([\lfp_{X_1} ~\Phi],...,[\lfp_{X_n} ~\Phi]).
	\]
	This formula does not depend on $X_1,...,X_n$, hence$
	(\calM,G),\theta \modelsa \beta([\lfp_{X_1} ~\Phi],...,[\lfp_{X_n} ~\Phi]).
	$\\
	$``\leftarrow"$:  For the other direction let $(\calM,G)$ be a model abstraction and $\theta$ an environment such that
	\[
	(\calM,G),\theta \modelsa \beta([\lfp_{X_1} ~\Phi],...,[\lfp_{X_n} ~\Phi]).
	\]
	Let $S_i \in V_{k_i}$ be the abstract relations defined by $[\lfp_{X_i} ~\Phi]$ for $i \in \{1,...,n\}$. $(S_1,...,S_n)$ is a fixed point of $F_{\Phi}^{\#}$, therefore by the definition of $F_{\Phi}^{\#}$ we have 
	\[
	\alpha_{k_i} \circ (F_{\Phi})_i(\gamma_{k_1}(S_1),...,\gamma_{k_n}(S_n)) \sqsubseteq S_i, \quad \text{for every}~ i \in \{1,...,n\},
	\]
	and as $\alpha_k,\gamma_k$ form a Galois connection for all $k$ we get 
	\[
	(F_{\Phi})_i(\gamma_{k_1}(S_1),...,\gamma_{k_n}(S_n)) \subseteq \gamma_{k_i}(S_i), \quad \text{for every}~ i \in \{1,...,n\}.
	\]
	Thus $(\gamma_{k_1}(S_1),...,\gamma_{k_n}(S_n))$ is a fixed point of $F_{\Phi}$ and by defining $\theta' := \theta[X_1 := \gamma_{k_1}(S_1),...,X_n := \gamma_{k_n}(S_n)]$, we obtain
	\[
	(\calM,G),\theta' \modelsa \bigwedge_{i=1}^n \forall \mybar{x_i} \left(\alpha_i(X_1,...,X_n,\mybar{x_i}) \impl X_i(\mybar{x_i})\right).
	\]
	By assumption $(\calM,G),\theta' \modelsa \beta(X_1,...,X_n)$ and therefore we conclude
	\[
	(\calM,G),\theta \modelsa \exists \mybar{X} \bigwedge_{i=1}^n \forall \mybar{x_i} \left(\alpha_i(X_1,...,X_n,\mybar{x_i}) \impl X_i(\mybar{x_i})\right) \wedge \beta(X_1,...,X_n).
	\]
	We get 2. directly from (\ref{abstactNonnengartproofPart2}), if we choose $S_1,...,S_n$  to be the abstract relations $R_1,...,R_n$.
\end{proof}

\end{document}
